\documentclass[11pt,a4paper]{article}

\usepackage[T1]{fontenc}
\usepackage[utf8]{inputenc}
\usepackage{libertinus}
\usepackage{microtype}
\usepackage[a4paper,margin=22mm,headheight=15pt]{geometry}
\usepackage{amsmath,amssymb,mathtools}
\usepackage{booktabs,longtable,tabularx,array}
\usepackage{enumitem}
\usepackage[most]{tcolorbox}
\usepackage{xcolor}
\usepackage{fancyhdr}
\usepackage[numbers,sort&compress]{natbib}
\usepackage{hyperref}

\definecolor{navy}{HTML}{163A5F}
\definecolor{teal}{HTML}{087E8B}
\definecolor{pale}{HTML}{EEF5F7}
\definecolor{warm}{HTML}{FFF6E8}
\definecolor{muted}{HTML}{4B5563}
\hypersetup{colorlinks=true,linkcolor=navy,citecolor=teal,urlcolor=teal,pdfauthor={DMD brain-state analysis planning},pdftitle={DMD and dynamical-system tracking: literature survey}}
\pagestyle{fancy}
\fancyhf{}
\lhead{DMD tracking literature survey}
\rhead{16 July 2026}
\cfoot{\thepage}
\setlength{\parindent}{0pt}
\setlength{\parskip}{0.55em}
\setlist{nosep,leftmargin=1.6em}
\renewcommand{\arraystretch}{1.18}
\newcolumntype{Y}{>{\raggedright\arraybackslash}X}
\newcolumntype{P}[1]{>{\raggedright\arraybackslash}p{#1}}

\newtcolorbox{keybox}[1][]{breakable,colback=pale,colframe=teal,boxrule=0.7pt,arc=2pt,left=7pt,right=7pt,top=6pt,bottom=6pt,title=#1,fonttitle=\bfseries}
\newtcolorbox{warnbox}[1][]{breakable,colback=warm,colframe=orange!70!black,boxrule=0.7pt,arc=2pt,left=7pt,right=7pt,top=6pt,bottom=6pt,title=#1,fonttitle=\bfseries}

\begin{document}
\hypersetup{pageanchor=false}

\begin{titlepage}
  \centering
  \vspace*{2.0cm}
  {\color{navy}\rule{\textwidth}{1.2pt}}\\[1.1cm]
  {\Huge\bfseries Dynamic Mode Decomposition and\\[0.25em]
  Dynamical-System Tracking\par}
  \vspace{0.55cm}
  {\Large A critical literature survey for single-cell calcium imaging across wake, NREM sleep, and isoflurane anaesthesia\par}
  \vspace{1.0cm}
  {\color{navy}\rule{\textwidth}{0.6pt}}\\[1.0cm]
  \begin{tcolorbox}[colback=pale,colframe=teal,width=0.92\textwidth,arc=2pt]
  \centering
  \textbf{Planning deliverable only.} No new DMD analysis, statistical test, or scientific result is implemented or claimed in this document.
  \end{tcolorbox}
  \vfill
  \begin{tabular}{rl}
    Prepared for: & \texttt{ca\_imaging\_network\_analysis}\\
    Source audit date: & 16 July 2026\\
    Companion document: & \texttt{02\_dmd\_brain\_state\_analysis\_plan.pdf}
  \end{tabular}
  \vfill
\end{titlepage}

\hypersetup{pageanchor=true}
\pagenumbering{roman}
\tableofcontents
\clearpage
\pagenumbering{arabic}

\section{Executive conclusions}

\begin{keybox}[The central correction]
Raut et al. did \emph{not} average spikes and then apply DMD to that population mean. Their Neuropixels preprocessing produced one scalar mean-rate observable for a separate arousal-manifold analysis. Their DMD was fitted to a multivariate widefield-calcium movie, with one observable per cortical pixel \citep{raut2025}. For the present data, the corresponding DMD state is therefore a \emph{neuron-by-time rate-like matrix}; averaging over neurons before DMD would destroy the spatial modes and reduce the model to a scalar autoregression.
\end{keybox}

The literature supports a two-level strategy:

\begin{enumerate}
  \item compare each window's full low-rank operator with a joint spectral--subspace distance; and
  \item track only bootstrap-stable conjugate pairs or invariant subspaces through time, allowing modes to appear, disappear, merge, and cross.
\end{enumerate}

Germain et al.'s spectral--Grassmann optimal-transport (SGOT) construction is well suited to the first level, but it is a \emph{pairwise metric}, not a temporal tracker \citep{germain2026}. Its optimal-transport coupling can seed adjacent-window correspondences. A true tracker still needs temporal consistency, birth/death penalties, uncertainty, and global optimization. Across animals, raw single-neuron mode subspaces cannot be compared directly because neuron identities and ambient dimensions differ. Cross-animal inference should therefore use scalar operator summaries, within-animal contrasts, or a preregistered common spatial representation.

The strongest confirmatory design is paired and hierarchical: awake versus NREM within sleep recordings, and awake versus anaesthesia within anaesthesia recordings. A direct three-state classifier is confounded by animal/session composition, circadian phase, and induction protocol. Overlapping windows are useful for visualization but are not independent replicates; inference must be at bout, session, and ultimately mouse level.

Finally, sliding estimators can create attractive but spurious trajectories in stationary autocorrelated data \citep{leonardi2015,hindriks2016}. Simulations, autocorrelation-preserving surrogates, rate-matched controls, grouped validation, and explicit uncertainty are therefore part of the proposed method rather than optional checks.

\section{Scope and review method}

This is a structured, critical narrative review, not a PRISMA-style systematic review or meta-analysis. It was designed to answer five methodological questions:

\begin{enumerate}
  \item What preprocessing and DMD procedure did Raut et al. actually use?
  \item What mathematical object does Germain et al.'s SGOT compare, and what assumptions does it impose?
  \item Which DMD estimators are credible under noisy, stochastic neural observations?
  \item Which methods address time variation, subspace/mode matching, and discrete brain-state dynamics?
  \item Which controls prevent calcium kinetics, rate changes, window overlap, or cross-animal geometry from masquerading as a change in dynamics?
\end{enumerate}

The two requested local PDFs were read in full, including Methods and appendices. The local Cell Reports paper describing this dataset was also audited \citep{kiyooka2026}. The public code and reporting details linked to the Raut and Germain papers were checked to resolve parameters absent from the prose.\footnote{\href{https://github.com/ryraut/arousal_dynamics}{Raut repository}, commit \texttt{dc7f7e24786a}; \href{https://github.com/thibaut-germain/SGOT}{Germain repository}, commit \texttt{09704a7f7ef5}.} The broader evidence map was assembled by citation chaining from these sources and targeted searches of primary, peer-reviewed articles on DMD, stochastic/noisy DMD, online and non-stationary system identification, Grassmann/operator metrics, neural DMD, switching state-space models, calcium deconvolution, and dynamic-connectivity null models. Preprints are identified as such. The cutoff was 16 July 2026.

Selection emphasized papers that contributed an estimator, validation result, neural application, or directly transferable failure analysis. General reviews were not used as substitutes for primary methods papers. This approach is rigorous about the decision at hand, but it does not estimate publication bias or claim exhaustive coverage.

\section{What Raut et al. actually did}

Local PDF locators are Methods ``Dataset 1/2'' preprocessing (p.~9), Methods ``Dynamic mode decomposition'' (pp.~11--12), the Figure~2c result (p.~4), and the Reporting Summary software list (p.~25). The separate official supplement contains no additional DMD parameter or validation result.

\subsection{Two pipelines that must not be conflated}

Raut et al. studied whether a delay embedding of a scalar arousal proxy could reconstruct multivariate brain dynamics \citep{raut2025}. Their spike and DMD procedures served different parts of that argument.

\begin{longtable}{@{}P{0.19\textwidth}P{0.28\textwidth}P{0.23\textwidth}P{0.22\textwidth}@{}}
\toprule
Data and role & Preprocessing immediately relevant here & Analysis target & Consequence for this project\\
\midrule
\endfirsthead
\toprule
Data and role & Preprocessing immediately relevant here & Analysis target & Consequence for this project\\
\midrule
\endhead
Allen Neuropixels; one arousal observable & Kilosort2 spikes; 30-min spontaneous epoch; quality-controlled neocortical units; 1.2-s bins; mean over units; outlier mitigation; alignment/resampling and 0.01--0.2-Hz filtering & One scalar population trace used alongside pupil, running, facial motion, LFP power, and other observables & Useful as a population-rate baseline/covariate. It is not the input on which their DMD modes were estimated.\\
\addlinespace
Widefield jRGECO1a; spatial brain dynamics & 128$\times$128 grid; atlas registration; detrending; spatial Gaussian smoothing; pixelwise $\Delta F/F$; haemodynamic correction; 0.01--0.2-Hz band-pass; brain mask; temporal mean-centering & A high-dimensional pixel-by-time matrix fitted with rank-2 randomized DMD & The relevant analogy is one neuronal observable per row, not the across-neuron mean. Their narrow band and rank 2 reflect a specific infraslow-wave question, not universal DMD defaults.\\
\bottomrule
\end{longtable}

The public Neuropixels code resamples 1.2-s \emph{counts} without visibly dividing by bin duration before averaging. Subsequent standardization makes this scale immaterial for their model, but it reinforces that the procedure should not be copied as a calibrated firing-rate estimator. Interpolation to 20 Hz also does not create information above the original bin rate. The code does not visibly reproduce every prose step (notably the neocortical-region restriction and median filtering), so the manuscript procedure and incidental code behaviour should be reported separately.

\subsection{Widefield preprocessing before DMD}

The reported sequence was spatial normalization; downsampling to 128$\times$128; co-registration to the Paxinos atlas; fifth-order polynomial detrending; 5$\times$5-pixel Gaussian spatial smoothing (1.3-pixel standard deviation); pixelwise $\Delta F/F$ relative to the run mean; haemoglobin-absorption correction; and 0.01--0.2-Hz band-pass filtering \citep{raut2025}. The final arrays were treated as 20-Hz signals. The public implementation retained 10,400 brain-mask pixels, temporally mean-centred each pixel without variance scaling, and removed edge segments associated with filtering and delay-model construction.

These choices isolate a 5--100-s infraslow process in ten-minute awake recordings. They do not establish that the same band is optimal for 39-s state bouts or that a 20-Hz interpolated signal has 20-Hz information. A window intended to estimate 0.01-Hz dynamics would require at least roughly 200--300 s to contain two or three cycles; the current dataset's minimum accepted state bout is only 300 frames (39.2 s). Window length and interpretable frequency floor must therefore be coupled explicitly.

\subsection{DMD formulation and exact implementation}

For state snapshots $Y=[y(t_1),\ldots,y(t_p)]$, Raut et al. formed adjacent matrices $Y^{(1)}$ and $Y^{(2)}$ and fitted
\begin{equation}
  Y^{(2)} \approx A Y^{(1)},
\end{equation}
with a low-rank representation
\begin{equation}
  y(t) \approx \sum_{k=1}^{r} b_k v_k e^{\omega_k t}.
\end{equation}
Their code used PyDMD's randomized DMD with \texttt{svd\_rank=2}, seed 1, oversampling 5, and one power iteration. The input orientation was 10,400 cortical pixels by time; a single global fit used all adjacent samples. It was not sliding-window DMD, Hankel DMD, forward--backward DMD, total-least-squares DMD, or optimized DMD.

They fitted DMD separately to the true held-out mouse's widefield series, the nonlinear pupil-delay reconstruction, and a no-embedding scalar reconstruction. The empirical feature was the spatial phase map of the first returned complex mode (not a mode explicitly ranked by amplitude or energy, because eigenvalue sorting remained off),
\begin{equation}
  \theta_i=\arg(v_{i1}),
\end{equation}
compared by an absolute circular correlation to remove arbitrary complex phase/sign. They did not report state differences, held-out DMD forecasts, eigenvalue distributions, decay rates, rank sensitivity, or window tracking. Thus the reusable precedents are snapshot construction, complex-mode interpretation, and phase-invariant comparison---not population averaging, rank 2, or the 0.01--0.2-Hz band.

\subsection{Translation to the present calcium data}

The present repository already stores OASIS-derived \texttt{spike\_deconv} and a 15-frame (1.96-s) Gaussian-smoothed \texttt{spike\_smoothed} signal \citep{friedrich2017,kiyooka2026}. These are event-rate proxies, not calibrated spikes per second. The faithful multivariate state is
\begin{equation}
  x(t)=\big[r_1(t),\ldots,r_N(t)\big]^\top,
\end{equation}
where the same within-session neurons remain in the same rows for every compared window. The population mean $N^{-1}\sum_i r_i(t)$ should be retained as a simple baseline and nuisance descriptor. A 1.2-s binned \texttt{spike\_deconv} arm can test sensitivity to the Raut timescale, but it should not be upsampled and treated as independent high-rate data.

\section{What Germain et al. contribute}

The local filename's ``2025'' refers to the arXiv preprint; the audited 37-page PDF is the accepted ICLR 2026 paper. Key locators are the assumptions (Section~3, p.~4), spectral measure and SGOT metric (Eqs.~4--5, p.~5), barycenter (Eq.~7, p.~6), and Grassmann appendix (pp.~18--20).

\subsection{A distribution over spectral atoms}

Germain et al. represent a finite-rank, non-defective Koopman/transfer generator through eigenvalues and spectral-projector subspaces \citep{germain2026}. If
\begin{equation}
  L=\sum_j \lambda_j f_j\otimes g_j,
\end{equation}
then each atom carries generator eigenvalue $\lambda_j$ and an eigenspace $\mathcal V_j$ in a Hilbert--Schmidt operator space. With multiplicity $m_j$,
\begin{equation}
  \mu(T)=\sum_j \frac{m_j}{m_{\mathrm{tot}}}\,
  \delta_{(\lambda_j,\mathcal V_j)}.
\end{equation}
For subspaces $U$ and $V$, their projector metric is
\begin{equation}
  d_G(U,V)=\|\Pi_U-\Pi_V\|_{\mathrm{HS}}.
\end{equation}
For equal dimension $q$ and principal angles $\theta_a$, this obeys
\begin{equation}
  d_G^2(U,V)=2\sum_{a=1}^{q}\sin^2\theta_a.
\end{equation}

The ground cost joins generator and geometric discrepancies,
\begin{equation}
 d_\eta\big((\lambda,U),(\lambda',V)\big)
 =\eta\,d_\lambda(\lambda,\lambda')+(1-\eta)d_G(U,V),
\end{equation}
and SGOT is the Wasserstein optimal-transport cost between the two spectral-atom distributions. This makes the comparison invariant to arbitrary mode ordering and capable of comparing unequal ranks or multiplicities.

\subsection{Assumptions that matter here}

The theory assumes time-homogeneous systems, isolated leading generator eigenvalues, a common observable/RKHS space, and non-defective finite-rank operators. Its finite-sample result is for regularized reduced-rank regression under equilibrium and regularity conditions; it does not directly cover ordinary DMD fitted to strongly dependent, overlapping calcium windows.

Three distinctions are critical:

\begin{enumerate}
  \item The full spectral projector depends on both right and left eigenfunctions. Comparing only right DMD modes on a Grassmann manifold is useful, but it is an approximation and should be named ``SGOT-inspired,'' not exact SGOT.
  \item The theoretical atom mass is eigenvalue multiplicity, not modal amplitude or biological energy. Energy-weighted transport is a reasonable sensitivity analysis but changes the metric's meaning.
  \item A common ambient space exists naturally within a recording, where the ROI identities are fixed. It does not exist across mice with different neurons and dimensions unless the observations are first mapped into a common parcellation, spatial basis, or learned representation.
\end{enumerate}

For real data, grouping a complex-conjugate pair into one real two-dimensional invariant subspace is often the stable object for tracking. Exact SGOT nevertheless retains the two distinct conjugate generator eigenvalues as separate atoms; a one-atom, positive-frequency pair representation is a useful but explicitly modified quotient construction.

\subsection{Why SGOT is not yet a tracker}

The transport plan provides a pairwise alignment. It does not assign persistent identities across three or more windows, penalize track acceleration, or model birth and death. Balanced OT also forces every unit of mass to match something. A temporal extension therefore requires dummy birth/death nodes, partial or unbalanced transport \citep{chizat2018}, or a global min-cost-flow graph across the entire bout. Near repeated eigenvalues and crossings, individual eigenvectors are unstable; the identifiable object is the joint invariant subspace.

The paper's Fr\'echet barycenter is relevant for state prototypes, but optimization is non-convex. State medoids---observed windows minimizing within-state SGOT distance---are a safer primary prototype; barycenters can remain exploratory. The public code uses continuous-time logarithms, balanced Earth-mover transport, order $p=2$, and uniform mass for simple atoms. Its rank-one projector scaling is not numerically identical to the paper's Hilbert--Schmidt convention, and the reported versus coded classification regularization differs. Reimplementation must select one documented convention and unit-test it rather than silently mixing paper and code normalizations.

\section{Evidence map beyond the two focal papers}

\subsection{Estimator choice under noise and stochasticity}

Classical exact DMD estimates a best-fit linear propagator from adjacent snapshots \citep{schmid2010,tu2014}. In neural observations, ordinary least squares is asymmetrically biased because both snapshot matrices contain noise. Total-least-squares, forward--backward, and noise-corrected variants address this failure \citep{dawson2016,hemati2017}. Subspace DMD targets stochastic Koopman analysis \citep{takeishi2017}; optimized DMD treats the problem as variable projection \citep{askham2018}; bagged optimized DMD supplies empirical spectral and modal uncertainty \citep{sashidhar2022}. These methods motivate a state-blind benchmark rather than declaring one estimator a priori solely because it appears in Raut.

\subsection{Time-varying systems and regime discovery}

Online/sliding DMD updates local operators for time-varying systems \citep{zhang2019online}. Multiresolution DMD separates slow and fast temporal components hierarchically \citep{kutz2016mrdmd}; windowed DMD with spectral clustering exposes local spectral regimes but also demonstrates the time--frequency trade-off \citep{dylewsky2019}. Non-stationary DMD estimates smoothly modulated modes and frequencies \citep{ferre2023}, whereas factorially switching DMD and recurrent switching linear dynamical systems represent recurring discrete regimes \citep{takeishi2018switching,linderman2017}. Low-rank time-varying autoregression imposes smooth shared structure across a sequence of local system matrices \citep{harris2021}.

These are competing scientific hypotheses, not merely algorithms. Sliding DMD assumes approximate stationarity within each window. Non-stationary DMD assumes gradual drift. Switching models assume a finite recurrent repertoire. Comparing their held-out predictions and transition recovery is more informative than presenting only a single smooth trajectory.

\subsection{Grassmann and operator-comparison precedents}

Principal angles provide a phase- and basis-invariant comparison between subspaces. Neural applications have already used DMD-mode subspaces and Grassmann projection kernels for decoding \citep{shiraishi2020,fukuma2024}. Dynamical similarity analysis compares fitted linear dynamics after alignment and is explicitly designed for neural population comparisons \citep{ostrow2023}. Operator metrics based on Perron--Frobenius/Koopman representations provide another coordinate-aware benchmark \citep{ishikawa2018,zhang2025koopstd}. Grassmann state-space regularizers can represent constant-position or constant-velocity subspace motion \citep{saadfalcon2024}, but smoothing must not be allowed to erase a true abrupt state transition.

\subsection{DMD in neuroscience}

DMD has exposed spatial--spectral patterns in ECoG, including sleep-spindle networks \citep{brunton2016neural}, and time-resolved resting-state networks in fMRI \citep{kunertgraf2019}. The latter used overlapping windows and simulations that showed rapid state changes are temporally smeared and can be undercounted. DMD has also been used on resting/task fMRI \citep{casorso2019} and with control to separate intrinsic dynamics from behaviour-transition inputs in whole-brain calcium imaging \citep{fieseler2020}. These studies establish feasibility but do not eliminate the need for calcium-specific observation controls or animal-level validation.

Brain-state research supplies complementary endpoints. Hidden-state analyses quantify occupancy, lifetime, and transition graphs across wake and NREM \citep{stevner2019}; anaesthesia recovery can proceed through metastable states rather than a monotonic path \citep{hudson2014}. Calcium slow waves alone can strongly separate wake, sleep, and anaesthesia \citep{brier2019,tortcolet2023}. Therefore, any claim that DMD captures a distinctive dynamical signature must show incremental information beyond mean rate, global components, spectral power, and autocorrelation.

\subsection{Calcium observation model and behavioural confounding}

OASIS infers sparse events under an explicit calcium-decay observation model \citep{friedrich2017}; smoothing and indicator kinetics alter apparent timescales, phase, and rotations. Simultaneous calcium/electrophysiology comparisons show that population dynamics can differ between the modalities \citep{wei2020}, and calcium-specific latent-dynamics methods are an important sensitivity benchmark \citep{koh2023}. Data-driven rate-bandwidth methods offer a principled alternative to arbitrary binning \citep{shimazaki2010}.

Movement and arousal can explain a large fraction of cortical activity \citep{stringer2019,musall2019}. This dataset lacks all of the continuous pupil, locomotion, respiration, and whisking covariates used by Raut, so state interpretation must acknowledge residual behavioural confounding. Within-session pairing and a quiet-wake sensitivity analysis reduce, but do not remove, this limitation.

\section{Methodological decision matrix}

\begin{longtable}{@{}P{0.18\textwidth}P{0.25\textwidth}P{0.27\textwidth}P{0.22\textwidth}@{}}
\toprule
Method & What it identifies & Main strength & Primary limitation / assigned role\\
\midrule
\endfirsthead
\toprule
Method & What it identifies & Main strength & Primary limitation / assigned role\\
\midrule
\endhead
Exact/randomized DMD & Local linear modes and eigenvalues & Transparent; direct Raut replication & Noise bias and no uncertainty; benchmark only\\
TLS / forward--backward / subspace DMD & Noise-robust local operator & Better matched to noisy stochastic observations & Variant-specific assumptions; estimator benchmark\\
BOP-DMD & Optimized spectra and bootstrap uncertainty & Natural stability filter and uncertainty & Greater computational cost; candidate primary if stable\\
Windowed DMD & Time-local operator sequence & Simple and interpretable trajectory & Window mixing, rank flicker, correlated estimates; main framework with strong nulls\\
Non-stationary DMD / TVAR & Smoothly changing dynamics & Borrows strength through time & Can blur abrupt transitions; competing smooth-drift model\\
HMM / rSLDS / switching DMD & Recurrent discrete regimes & Direct occupancy/dwell/transition model & Model-order and identifiability sensitivity; competing metastable-state model\\
Right-mode Grassmann distance & Spatial/invariant-subspace rotation & Simple, phase invariant & Ignores eigenvalues and left modes; subspace diagnostic\\
Exact SGOT & Generator spectra plus spectral projectors & Permutation invariant; ranks may differ & Needs common space and left/right spectral information; primary operator distance within session\\
Unbalanced OT / min-cost flow & Persistent mode lineages & Handles birth/death and global identity & Proposed extension, not Germain's original method; secondary track analysis\\
Dynamical similarity analysis & Aligned latent dynamics across datasets & Useful cross-animal benchmark & Depends on learned representation/alignment; exploratory\\
\bottomrule
\end{longtable}

\section{Recommended synthesis}

The literature supports the following architecture, developed into an executable specification in the companion plan:

\begin{enumerate}
  \item Use the full per-neuron smoothed event-proxy matrix as the primary DMD state, with binned deconvolved events and $\Delta F/F$ as sensitivity arms. Preserve mean-rate and sparsity summaries rather than erasing them through independent windowwise standardization.
  \item Build windows only inside contiguous acquisition and state bouts. Use a fixed state-blind rank within each analysis configuration, and choose estimator/rank by grouped held-out prediction plus bootstrap subspace stability.
  \item Represent complex-conjugate pairs as real two-dimensional invariant subspaces. Use the continuous-time generator $\omega=\log(\lambda)/\Delta t$ and explicitly guard against Nyquist/complex-log ambiguity.
  \item Report three complementary trajectories: a dominant DMD subspace trajectory; an SGOT operator-distance trajectory; and high-confidence individual/grouped mode lineages from a global birth/death-aware tracker.
  \item Compare states with within-mouse paired effects and leave-one-mouse-out validation. Do not split overlapping windows randomly or use window count as sample size.
  \item Require simulation recovery, stationary-surrogate calibration, prediction outperformance, and incremental value beyond rate/global-power baselines before interpreting a state difference as dynamical.
\end{enumerate}

\begin{warnbox}[Claims the present evidence does not support]
The literature does not justify calling the deconvolved signal an absolute firing rate; copying Raut's population average or 0.01--0.2-Hz filter; interpreting every DMD eigenvector as an identifiable biological mode; treating pairwise SGOT couplings as persistent tracks; comparing raw neuronal subspaces across mice; treating overlapping windows as replicates; or interpreting sleep--anaesthesia similarity as a causal signature of consciousness without addressing protocol and circadian confounds.
\end{warnbox}

\section{Evidence gaps and reporting requirements}

No reviewed paper validates the complete combination of thousands of deconvolved calcium traces, short state-homogeneous windows, nonnormal DMD operators, SGOT geometry, temporal births/deaths, and mouse-level sleep/anaesthesia inference. The proposed analysis is therefore method development as well as application. Every departure from published SGOT---energy weights, right-mode-only geometry, unbalanced transport, temporal smoothness---must be named and reported separately from the exact metric.

The eventual report should expose the full multiverse of window length, smoothing/binning, rank, estimator, lag, spectral scaling, and SGOT weight; publish per-mouse effects; show stationary null trajectories beside observed trajectories; and provide failure diagnostics for eigengaps, conditioning, bootstrap reproducibility, and assignment confidence. Negative or unstable results would still answer the scientific question by bounding which dynamical differences the available observation duration and signal quality can resolve.

\clearpage
\sloppy
\bibliographystyle{unsrtnat}
\bibliography{dmd_tracking_references}

\end{document}
